\documentclass{Latex/Classes/PhDthesisPSnPDF}
\usepackage[utf8]{inputenc}
\usepackage[T1]{fontenc}
\usepackage{geometry}
\usepackage{graphicx}
\usepackage{float}
\usepackage{hyperref}
\usepackage[english,indonesian]{babel}

\title{Optimasi Pemrosesan Template LaTeX Menggunakan Pendekatan Berbasis AI}
\author{Academic Support Agent}
\date{\today}

\begin{document}
\maketitle

\begin{abstract}

\end{abstract}

\chapter{PENDAHULUAN}


\chapter{TINJAUAN PUSTAKA}


% Original content (commented out):
% % -----------------
%  UNIVERSITY OF LIMERICK

% Template for PhD Thesis by Publication
% Prepared by Hafiz Ahmad Awais Chaudhary
% https://www.linkedin.com/in/ahmad91/
% 2023
% Department of Computer Science and Information Systems - University of Limerick, Ireland
 

% Based on Giuseppe Torre, Harish Bhandari's and Jakob Suckale templates  
% ------------------

%: Style file for Latex
% Most style definitions are in the external file PhDthesisPSnPDF.
% In this template package, it can be found in ./Latex/Classes/
\documentclass[openright,12pt]{Latex/Classes/PhDthesisPSnPDF}
%%%%%openany book twoside
\usepackage{natbib}
\usepackage{url}
\usepackage{alltt} 
\usepackage{multirow}
\usepackage{color}
\usepackage{graphicx}
\usepackage{subfig}
\usepackage{pdfpages}
\usepackage{float}
\usepackage{array}
\usepackage{tcolorbox}
\usepackage[acronym,toc]{glossaries}

\usepackage{xcolor}
%\usepackage{geometry}
\def\UrlBreaks{\do\/\do-}

\setcitestyle{square}

\makeglossaries

\newcommand{\urlwofont}[1]
{
\urlstyle{same}\url{#1}
}

% to display grey publication box
\definecolor{lightgrey}{rgb}{0.9, 0.9, 0.9}

% to display Math
\usepackage{mathtools}

% to Display C++ code
\definecolor{c++red}{rgb}{0.6,0,0} % for strings
\definecolor{c++green}{rgb}{0.25,0.5,0.35} % comments
\definecolor{c++purple}{rgb}{0.5,0,0.35} % keyword
\usepackage{listings}
\lstset{language=C++}


%% Define a new 'Leo' style for the package that will use a smaller font.
\makeatletter  
\def\url@leostyle{%
  \@ifundefined{selectfont}{\def\UrlFont{\sf}}{\def\UrlFont{\small\ttfamily}}}
\makeatother
%% Now actually use the newly defined style.
\urlstyle{leo}





%: ----------------------------------------------------------------------
%:                  TITLE PAGE: name, degree,..t
% ----------------------------------------------------------------------
% below is to generate the title page with crest and author name

% If output to PDF then put the following in the PDF header
\ifpdf  
    \pdfinfo { /Title  (PhD and MPhil Thesis Classes)
               /Creator (TeX)
               /Producer (pdfTeX)
               /Author (Name Surname email@email.com)
               /CreationDate (D:YYYYMMDDhhmmss)  %format D:YYYYMMDDhhmmss
               /ModDate (D:YYYYMMDDhhmm)
               /Subject (Main subjects of your thesis)
               /Keywords (put here few keywords relevant to your PhD work) }
    \pdfcatalog { /PageMode (/UseOutlines)
                  /OpenAction (fitbh)  }
\fi


\title{Title}



% ----------------------------------------------------------------------
% The section below defines www links/email for authors and institutions
% They will appear on the title page of the PDF and can be clicked
\ifpdf
  \author{\href{mailto:email@ul.ie}{\textbf{ student name}}}
%  \cityofbirth{born in XYZ} % uncomment this if your university requires this
%  % If the city of birth is required, also uncomment 2 sections in PhDthesisPSnPDF
%  % Just search for the "city" and you'll find them.

% Here the links to the Institution you belong to
  \researchcentre{\href{https://lab.ie/}{lab name}}
  \faculty{\href{https://www.ul.ie/scieng/schools-and-departments/department-computer-science-and-information-systems}{Department of Computer Science and Information Systems}}
  \department{\href{https://www.ul.ie/scieng}{Faculty of Science and Engineering}}
  \university{\href{https://www.ul.ie/}{University of Limerick, Ireland}}

  % The crest is a graphics file of the logo of your research institution.
  % Place it in ./0_frontmatter/figures and specify the width
  \crest{\includegraphics[scale=0.45]{logonew2020.jpg}}
  
% If you are not creating a PDF then use the following. The default is PDF.
%\else
 % \author{YourName}
%  \cityofbirth{born in XYZ}
 % \collegeordept{CollegeOrDept}
  %\university{University}
  %\crest{\includegraphics[width=4cm]{logo.png}}
%\fi

%\renewcommand{\submittedtext}{change the default text here if needed}

\degree{Philosophi\ae Doctor (PhD) 2023}
\degreedate{}


% ----------------------------------------------------------------------
       
% turn of those nasty overfull and underfull hboxes
\hbadness=10000
\hfuzz=50pt


%: --------------------------------------------------------------
%:                  FRONT MATTER: dedications, abstract,..
% --------------------------------------------------------------

\begin{document}

%: Macro file for Latex
% Macros help you summarise frequently repeated Latex commands.
% Here, they are placed in an external file /Latex/Macros/MacroFile1.tex
% A macro that you may use frequently is the figure macro (see introduction.tex)
\include{Latex/Macros/MacroFile1}
%\language{english}



% sets line spacing
\renewcommand\baselinestretch{1.2}
\baselineskip=18pt plus1pt


%: ----------------------- generate cover page ------------------------

\maketitle  % command to print the title page with the above variables

\newpage
%: ----------------------- cover page back side ------------------------
% Your research institution may require reviewer names, etc.
% This cover back side is required by 

\begin{table}
\vspace*{-1cm}
\hspace{0.5cm}
\renewcommand{\arraystretch}{1}
\begin{tabular}{p{2.2cm} p{10.5cm} }

\textbf{Supervisor:} & Prof.    \\
 & Chair of    \\
 & Department of Computer Science and Information Systems  \\
 & University \emph{of} Limerick  \\
 & Ireland \\
 
 &  \\
 
\multirow{2}{2.2cm}{\textbf{Internal Examiner:}} & Dr. ...  \\
 & Department of Computer Science and Information Systems  \\
 & University \emph{of} Limerick  \\
 & Ireland \\


 &  \\
 
\multirow{2}{2.2cm}{\textbf{External Examiner:}} & Prof. Dr....  \\
 & Division of Software ...s  \\
  & University  \\
  & Country  \\

 &  \\

\textbf{Chair:} & Dr. ...  \\
 & Department of Computer Science and Information Systems  \\
 & University \emph{of} Limerick  \\
 & Ireland \\
 
 &  \\

\multicolumn{2}{@{} l}{\textbf{\hspace{0.1cm} Day of the defence: ..... 2023}}\\
 &  \\
  &  \\
   &  \\
   &  \\
 
 \multicolumn{2}{@{} r}{\textbf{Signature from the head of PhD committee:}}\\
 
\end{tabular}

\end{table}




%: ----------------------- abstract ------------------------

% Your institution may have specific regulations if you need an abstract and where it is to be placed in the document. The default here is just after title.
%: Declaration of originality



%\cleardoublepage 
\include{0_frontmatter/declaration}
%\includepdf{0_frontmatter/sign2.pdf}

\include{0_frontmatter/acknowledgement}

\include{0_frontmatter/dedication}
\include{0_frontmatter/abstract}

%: ----------------------- contents ------------------------

\setcounter{secnumdepth}{2} % organisational level that receives a numbers
\setcounter{tocdepth}{3}    % print table of contents for level 3
\tableofcontents            % print the table of contents
% levels are: 0 - chapter, 1 - section, 2 - subsection, 3 - subsection


%: ----------------------- list of figures/tables ------------------------
\include{0_frontmatter/glossary}

\listoftables  % print list of tables
\listoffigures	% print list of figures

%\printglossary[type=\acronymtype]
\printglossary[type=\acronymtype,nonumberlist, toctitle=List of Abbreviations, title=List of Abbreviations]
%\printglossary[type=\acronymtype,nonumberlist]



%: ----------------------- glossary ------------------------

% Tie in external source file for definitions: /0_frontmatter/glossary.tex
% Glossary entries can also be defined in the main text. See glossary.tex
%\include{0_frontmatter/glossary} 
%
%\begin{multicols}{2} % \begin{multicols}{#columns}[header text][space]
%\begin{footnotesize} % scriptsize(7) < footnotesize(8) < small (9) < normal (10)
%
%\printnomenclature[1.5cm] % [] = distance between entry and description
%\label{nom} % target name for links to glossary
%
%\end{footnotesize}
%\end{multicols}



%: --------------------------------------------------------------
%:                  MAIN DOCUMENT SECTION
% --------------------------------------------------------------

% the main text starts here with the introduction, 1st chapter,...
\mainmatter

\renewcommand{\chaptername}{} % uncomment to print only "1" not "Chapter 1"
\include{0_frontmatter/presentation}


%: ----------------------- Part I ------------------------

%\addcontentsline{toc}{part}{Part 1:}
\part{Introduction}
\label{part:part1}

\include{1_introduction/introduction}
\include{Literature/Literature}
\include{methods/chapter}	 
\include{results/results} 
\include{conclusions/conclusions} 
%: ----------------------- Part II ------------------------
\part{The Papers}
\label{part:part2}

\include{paper_j_syrcose/paper_j_syrcose}
\include{paper_c_isola_interoperability/paper_c_isola_interoperability}
\include{paper_j_eceasst/paper_j_eceasst}
\include{paper_j_mdpi/paper_j_mdpi}
\include{paper_c_isola_dtsm/paper_c_isola_dtsm}

% the following are not attached with the thesis, so ignore this section
%%%%%%%%%%%%%%%%%%%%%%\include{paper_c_isola_industrial/paper_c_isola_industrial}
%%%%%%%%%%%%%%%%%%%%%%%%\include{paper_c_isola_dime/paper_c_isola_dime}
%%%%%%%%%%%%%%\include{paper_c_ifip_iot/paper_c_ifip_iot}




%: ----------------------- Part III ------------------------
\part{Appendixes}
\label{part:part3}
%\include{0_frontmatter/part3}
%\section{Poster}
\section*{Exemplary Code}
\addcontentsline{toc}{chapter}{Exemplary Code}
\label{appendix:code}


\newpage

\phantomsection
\addcontentsline{toc}{chapter}{PhD Poster}
if you want to include any pdf then give path here
%\includepdf{poster.pdf}
\label{poster_tech}
\section*{Project Funding}
\addcontentsline{toc}{chapter}{Project Funding}


This PhD is funded by the .



% for reference list see line 361
% description of lab methods
% AC comments
% here we start the Appendix
%\appendix

% Definition of each Chapter in the Appendix is done via the following two lines
%\chapter*{Appendix A}
%\section*{Insert a figure}

% here we open a new chapter in Appendix etc. etc..
%\chapter*{Appendix A}
%\section*{Title}

%Type here you text as usual . . .


%\addcontentsline{toc}{chapter}{Appendix A: .....}
%\chaptermark{Appendix}
%\markboth{Appendix}{Appendix}

% AC comments


% --------------------------------------------------------------
%:                  BACK MATTER: appendices, refs,..
% --------------------------------------------------------------

% the back matter: appendix and references close the thesis






%\chapter*{Appendix C}
%\section*{Code for estimating attitude}

%AHRS\_TRIAD.C




%\lstset{language=C++}
%\lstset{basicstyle=\tiny, 
%breaklines=true, 
%frame=shadowbox, 
%rulesepcolor=\color{c++green}, 
%keywordstyle=\color{c++purple}\bfseries, 
%stringstyle=\color{c++red}, 
%commentstyle=\color{c++green},}  % edit these parameters if you want to change the way your code is displayed

%Load your code file below
%\lstinputlisting[fontadjust=\true]{7_appendixes/ahrs_triad_mote.c}

%\addcontentsline{toc}{chapter}{Appendix C: Insert C++/Java et al. Code in Appendix}
%\chaptermark{Appendix}
%\markboth{Appendix}{Appendix}

%





%: ----------------------- bibliography ------------------------


%\begin{small} % tiny(5) < scriptsize(7) < footnotesize(8) < small (9)
%\bibliographystyle{Latex/IEEEtran_bib}
%\bibliographystyle{Latex/splncs04}
%\renewcommand{\bibname}{References} % changes the header; default: Bibliography
%\bibliography{References-Bibtex} % adjust this to fit your BibTex file
%\end{small}

%\end{multicols}

% --------------------------------------------------------------
% Various bibliography styles exit. Replace above style as desired.

% in-text refs: (1) (1; 2)
% ref list: alphabetical; author(s) in small caps; initials last name; page(s)
%\bibliographystyle{Latex/Classes/PhDbiblio-case} % title forced lower case
%\bibliographystyle{Latex/Classes/PhDbiblio-bold} % title as in bibtex but bold
%\bibliographystyle{Latex/Classes/PhDbiblio-url} % bold + www link if provided

%\bibliographystyle{Latex/Classes/jmb} % calls style file jmb.bst
% in-text refs: author (year) without brackets
% ref list: alphabetical; author(s) in normal font; last name, initials; page(s)

%\bibliographystyle{plainnat} % calls style file plainnat.bst
% in-text refs: author (year) without brackets
% (this works with package natbib)

%%%%%%%%%\chapter*{References}
%%%%%%%%%%%\addcontentsline{toc}{chapter}{References}

% -------------------------------------------------------------
\Urlmuskip=0mu plus 1mu

\bibliographystyle{Latex/IEEEtran_bib}
\bibliography{References-Bibtex}
\end{document}


\end{document}